\begin{frame}{같은 데이터, 다른 운명: 로지스틱회귀 vs 2층 신경망}
  \small
  같은 XOR 4개, 같은 경사하강법($\alpha = 1.0$, 5000 에폭):
  \vskip0.5em
  \begin{tabular}{@{}lcc@{}}
    \toprule
    모델 & 최종 예측 & 손실 \\
    \midrule
    로지스틱회귀 (은닉층 없음) & $[0.500, 0.500, 0.500, 0.500]$ & 0.125 \\
    2층 신경망 (은닉 4개) & $[0.016, 0.977, 0.976, 0.030]$ & 0.0003 \\
    \bottomrule
  \end{tabular}
  \vskip0.6em
  \begin{itemize}
    \item 로지스틱회귀: 네 점을 \textbf{전부 0.5(완전 불확신)}으로 예측한 채
      0.125에서 \textbf{멈춤} — 앞에서 본 부등식 모순이 수렴을 원천 차단
    \item 2층 신경망: 0.0003까지 내려 네 점을 정확히 분리
  \end{itemize}
  \begin{block}{구조의 차이가 운명을 정한다}
    ``은닉층이 왜 필요한가''를 가장 극적으로 보여주는 최소 예시.
  \end{block}
\end{frame}
